% paper.tex -- draft, A&A LaTeX class (aa.cls v7.0, EDP Sciences)
% Mixed poloidal-toroidal equilibria beyond barotropy.
% TODO: title, author list, affiliations, introduction, conclusions.
% TODO: the non-barotropic references were entered from memory -- check
% every volume and page against ADS before this leaves the machine.

\documentclass{aa}

\usepackage{graphicx}
\usepackage{amsmath}
\usepackage{amssymb}
\usepackage{txfonts}
\usepackage{natbib}

\begin{document}

   \title{Mixed poloidal-toroidal fields in massive white dwarfs
          beyond barotropy}

   \subtitle{Draft -- work in progress}

   \author{TBD\inst{1}
          }

   \institute{Affiliation TBD\\
              \email{tbd@example.org}
             }

   \date{Draft \today}

   \abstract
   % context
   {A white dwarf supported above the Chandrasekhar mass by a magnetic
    field needs most of that field to be toroidal, because the toroidal
    branch is the one that reaches the required energy. It also needs a
    poloidal component, or it has no exterior dipole and nothing for
    magnetic braking or an observer to act on.}
   % aims
   {We ask whether a self-consistent equilibrium can carry both at once,
    whether differential rotation then lifts it into the $1.8$--$2\,M_\odot$
    range of interest for ultramassive white dwarfs, and, if not, which
    assumption blocks it.}
   % methods
   {We build axisymmetric non-barotropic equilibria by integrating
    $\mathrm{d}P = g\,\mathrm{d}u$ along equipotentials and recovering the
    density from radial force balance, closing the system with a
    zero-temperature degenerate gas of variable composition so that the
    mean molecular weight per electron is diagnosed pointwise rather than
    imposed. Differential rotation enters through a $j$-constant law. We
    scan $908$ models and accept one only if it is simultaneously in
    equilibrium to a virial error below $10^{-3}$, below the Landau field,
    not shedding mass, of admissible composition on a mass-weighted
    pointwise test, and Ledoux-stable above the solver's own measured
    noise floor.}
   % results
   {The barotropic twisted torus saturates in this code at
    $E_{\mathrm{tor}}/E_{\mathrm{mag}} = 0.04$. Non-barotropy exceeds
    that, but the composition stratification required is Ledoux-unstable
    in all $126$ models where it is large enough to measure. Force balance,
    not composition, caps the toroidal fraction near
    $E_{\mathrm{tor}}/|W| \simeq 0.01$, so
    $B_\varphi/|\mathbf{B}_{\mathrm{p}}| \lesssim 0.6$ in every
    equilibrium, with or without rotation. Differential rotation reaches
    $2.85\,M_\odot$ in genuine equilibrium, but mass and composition trade
    monotonically and above $1.5\,M_\odot$ the solution demands
    $\mu_e < 2$, which no cold composition supplies. Exactly one model of
    $908$ passes every test: $1.278\,M_\odot$, no toroidal field, and
    lighter than its own field-free background.}
   % conclusions
   {Supplying the missing pressure thermally would require
    $4\times10^9$\,K at $1.4\,M_\odot$ and $2$--$5\times10^{10}$\,K at
    $2\,M_\odot$, against a pycnonuclear carbon ignition already active at
    these densities. An ultramassive white dwarf supported this way does
    not wait: the conditions it needs are ones under which carbon has
    already ignited.}

   \keywords{white dwarfs -- stars: magnetic field --
             magnetohydrodynamics (MHD) -- methods: numerical}

   \maketitle

%-------------------------------------------------------------------

\section{Introduction}

TBD.

%-------------------------------------------------------------------

\section{What the equilibrium forces, and what barotropy adds}
\label{sec:forced}

It is worth separating the constraints, because they are usually quoted
together and only one of them is negotiable.

Consider an axisymmetric star in magnetostatic equilibrium, with no
flows, and write the field through the flux function $u = \varpi
A_\varphi$,
\begin{equation}
\mathbf{B} = \frac{1}{\varpi}\,\nabla u \times
             \hat{\mathbf{e}}_\varphi
             + B_\varphi \hat{\mathbf{e}}_\varphi .
\label{eq:decomposition}
\end{equation}
Neither $\nabla P$ nor $\rho\nabla\Phi$ has an azimuthal component, so
the azimuthal component of the Lorentz force must vanish by itself. That
component is proportional to
$(\mathbf{B}_{\mathrm{p}}\cdot\nabla)(\varpi B_\varphi)$, so
\begin{equation}
\varpi B_\varphi = \beta(u),
\label{eq:beta}
\end{equation}
that is, $\varpi B_\varphi$ is constant along each poloidal field line.
Outside the star the field is in vacuum and $B_\varphi$ vanishes, so
$\beta$ must vanish on every line that reaches the surface. Writing
$u_s$ for the largest value of $u$ on the stellar surface, the toroidal
field is confined to the closed-line region $u > u_s$. This is the
twisted torus, and it is what our companion paper measures directly as a
separatrix.

The point to make is that \emph{none of this used barotropy}.
Equation~(\ref{eq:beta}) and the confinement follow from axisymmetry and
from the exterior being field-free. Generalising the equation of state
does not release them.

What barotropy does add is a restriction on the poloidal source. If
$\rho = \rho(P)$ then $\rho^{-1}\nabla P = \nabla h$ and the whole
left-hand side of the Euler equation is a gradient, so
$\rho^{-1}\mathbf{f}_L$ must be one too. Since the Lorentz force of
Eq.~(\ref{eq:decomposition}) is everywhere parallel to $\nabla u$, this
forces $\rho^{-1}\mathbf{f}_L = M'(u)\nabla u$ for some function $M$, and
the equilibrium collapses to the Grad--Shafranov equation
\citep{grad1958,shafranov1966}
\begin{equation}
\Delta^\ast u = -4\pi\varpi^2 \rho\, M'(u) - \beta(u)\beta'(u),
\label{eq:gs}
\end{equation}
with the Bernoulli integral $h + \Phi - M(u) = \mathrm{const}$. The
source of the poloidal field is pinned to the form $\rho M'(u)$: it must
be proportional to the density, with a coefficient constant on flux
surfaces.

That single restriction is the subject of this paper. Relaxing it is what
$\rho = \rho(P,s)$ buys, because then $\rho^{-1}\nabla P$ is no longer a
gradient, the baroclinic term $\nabla\rho\times\nabla P$ is available to
balance what the gradient cannot, and the poloidal source becomes free
\citep{mastrano2011,armaza2015}. It is also the route by which
\citet{ciolfi2013} reach large toroidal fractions in relativistic stars,
where the barotropic family had been limited to a few percent
\citep{lander2009,ciolfi2009,duez2010}.

For a white dwarf the physical case for it is straightforward: a real
white dwarf is stratified, in temperature and in composition, so
barotropy is the simplification and not the realism. The price is that
the equilibrium no longer determines itself -- a stratification has to be
specified, or derived and then checked for physical plausibility.

%-------------------------------------------------------------------

\section{The barotropic ceiling, measured}
\label{sec:ceiling}

Before building a solver to beat a ceiling, the ceiling should be
measured in the same code, on the same star. The value usually quoted for
the barotropic twisted torus is a few percent of magnetic energy in the
toroidal component; what follows is that number obtained here.

The background is a field-free configuration at
$\rho_c = 10^9$\,g\,cm$^{-3}$ and $\mu_e = 2$, giving
$M = 1.346\,M_\odot$, $R_{\mathrm{eq}} = 2.46\times10^8$\,cm and
$|W| = 2.257\times10^{51}$\,erg. On it we solve Eq.~(\ref{eq:gs}) with
$M'(u) = k_0$ constant, $k_0 = 10^{-13}$, and
\begin{equation}
\beta(u) = \beta_0
           \left(\frac{u-u_s}{u_{\mathrm{norm}}}\right)^{\zeta}
           \quad \text{for } u > u_s, \qquad 0 \text{ otherwise},
\label{eq:betaform}
\end{equation}
with $u_{\mathrm{norm}}$ fixed from the purely poloidal solution so that
the functional form does not drift as $u$ evolves. The equation is
nonlinear through $\beta\beta'(u)$ and is iterated to a relative change
below $10^{-8}$. We sweep $\beta_0$, quoted through the field it asks
for, $\beta_0 = B_{\mathrm{ref}}\varpi_{\mathrm{ref}}$ with
$\varpi_{\mathrm{ref}} = R_{\mathrm{eq}}/2$, and repeat for
$\zeta = 1$, $1.1$ and $2$.

At $\beta_0 = 0$ the poloidal solution has
$\max|\mathbf{B}_{\mathrm{p}}| = 2.45\times10^{12}$\,G, and its
closed-line region occupies $37.7\%$ of the stellar volume. Space is
therefore not what limits the toroidal component.

The sweep is summarised in Table~\ref{tab:ceiling}. The highest converged
toroidal fraction is
\begin{equation}
\left.\frac{E_{\mathrm{tor}}}{E_{\mathrm{mag}}}\right|_{\max}
  = 0.040,
\end{equation}
at $\zeta = 1$ and $B_{\mathrm{ref}} = 10^{12}$\,G, reached after $428$
iterations. Steeper $\beta$ profiles do worse: $\zeta = 1.1$ converges up
to $0.029$ and $\zeta = 2$ only to $0.003$. Above those values the
iteration leaves the basin entirely rather than degrading gracefully.

Two things follow. The first is that the few-percent figure quoted from
the literature is reproduced here, on a white dwarf rather than a neutron
star, which is worth knowing before it is used as an argument.

The second is quantitative and much sharper. At the ceiling the toroidal
field holds $E_{\mathrm{tor}}/|W| = 1\times10^{-5}$. A
$2\,M_\odot$ configuration of the kind our companion paper certifies
requires $E_{\mathrm{tor}}/|W| = 0.203$. The barotropic twisted torus is
not marginally short of what mass support needs; it is short by four
orders of magnitude. Whatever supports a $2\,M_\odot$ white dwarf, it is
not a self-consistent barotropic mixed equilibrium.

\begin{table}
\caption{Converged barotropic twisted-torus configurations. Rows above
the highest converged $B_{\mathrm{ref}}$ for each $\zeta$ are omitted
because the iteration diverges rather than returning a worse
equilibrium.}
\label{tab:ceiling}
\centering
\begin{tabular}{l l l l}
\hline\hline
$\zeta$ & $B_{\mathrm{ref}}$ (G) & $E_{\mathrm{tor}}/E_{\mathrm{mag}}$ &
$\max|B_\varphi|$ (G) \\
\hline
$1.0$ & $10^{11}$          & $0.004$ & $9.84\times10^{10}$ \\
$1.0$ & $3\times10^{11}$   & $0.029$ & $2.97\times10^{11}$ \\
$1.0$ & $10^{12}$          & $0.040$ & $7.23\times10^{11}$ \\
$1.1$ & $3\times10^{11}$   & $0.029$ & $3.10\times10^{11}$ \\
$2.0$ & $10^{11}$          & $0.003$ & $1.02\times10^{11}$ \\
\hline
\end{tabular}
\end{table}

\subsection{What this measurement does not settle}

Two caveats, both real.

The density is held fixed. The sweep solves Eq.~(\ref{eq:gs}) on the
field-free background rather than iterating the stellar structure
alongside it. The confinement of Sect.~\ref{sec:forced} is geometric and
survives that, and at these field strengths the back-reaction on the
structure is small, but the ceiling quoted here is a ceiling at fixed
density and should be called that.

The failure mode is a failure of the iteration. Above the values in
Table~\ref{tab:ceiling} an under-relaxed fixed-point scheme stops
converging. That is consistent with the equilibrium ceasing to exist, and
it is what the literature reports, but the two are not the same
statement, and separating them needs a continuation method that follows
the branch rather than a sweep that restarts at each point. Until that is
done, the honest claim is that the branch is unreachable here above
$4\%$, not that it is empty.

%-------------------------------------------------------------------

\section{Is there room for it? The confinement and the Landau limit}
\label{sec:room}

Two questions decide whether the construction is worth building: whether
it can give a toroidal field larger than the poloidal one, and whether it
keeps the mass high. The first is about the source, which non-barotropy
frees. The second is about the confinement, which it does not.

The difficulty is quantitative. The certified $2\,M_\odot$ configuration
carries $E_{\mathrm{tor}}/|W| = 0.203$ with $B_\varphi = K\rho\varpi$,
which is non-zero everywhere the density is. A mixed equilibrium must fit
the same energy into the closed-line region alone, and less volume at the
same energy means a larger peak field -- against a configuration that
already peaks at $4.28\times10^{13}$\,G, or $0.97\,B_c$ with
$B_c = m_e^2c^3/e\hbar = 4.414\times10^{13}$\,G. There is almost no
headroom before a field-independent equation of state stops describing
its own star.

We therefore ask, on the density of that configuration, what a confined
$\beta(u)$ of the form (\ref{eq:betaform}) with $\zeta = 1$ would cost.
Since $E_{\mathrm{tor}}$ is quadratic in $\beta_0$, the amplitude that
reproduces $E_{\mathrm{tor}}/|W| = 0.203$ follows by scaling, and the
question becomes what peak field it implies. Table~\ref{tab:room} gives
the answer over three decades of poloidal amplitude.

\begin{table*}
\caption{What a confined toroidal field costs, at fixed
$E_{\mathrm{tor}}/|W| = 0.203$, on the density of the certified
$2\,M_\odot$ configuration. Ratios $B_{\mathrm{t}}/B_{\mathrm{p}}$ are of
peak amplitudes, $E_{\mathrm{t}}/E_{\mathrm{p}}$ of energies.}
\label{tab:room}
\centering
\begin{tabular}{l l l l l l l l}
\hline\hline
$k_0$ & $B_{\mathrm{pole}}$ (G) & closed vol. &
$\max|B_\varphi|$ (G) & $/B_c$ & $\max|\mathbf{B}_{\mathrm{p}}|$ (G) &
$/B_c$ & $B_{\mathrm{t}}/B_{\mathrm{p}}$ \\
\hline
$10^{-13}$        & $5.9\times10^{10}$ & $37.2\%$ & $1.58\times10^{13}$ & $0.36$ & $1.63\times10^{12}$ & $0.04$ & $9.7$ \\
$3\times10^{-13}$ & $1.8\times10^{11}$ & $37.2\%$ & $1.58\times10^{13}$ & $0.36$ & $4.90\times10^{12}$ & $0.11$ & $3.2$ \\
$10^{-12}$        & $5.9\times10^{11}$ & $37.2\%$ & $1.58\times10^{13}$ & $0.36$ & $1.63\times10^{13}$ & $0.37$ & $1.0$ \\
$3\times10^{-12}$ & $1.8\times10^{12}$ & $37.2\%$ & $1.58\times10^{13}$ & $0.36$ & $4.90\times10^{13}$ & $1.11$ & $0.3$ \\
$10^{-11}$        & $5.9\times10^{12}$ & $37.2\%$ & $1.58\times10^{13}$ & $0.36$ & $1.63\times10^{14}$ & $3.70$ & $0.1$ \\
\hline
\end{tabular}
\end{table*}

Three results, and the first two are favourable.

The confinement is cheaper than the space argument suggests. Reproducing
$E_{\mathrm{tor}}/|W| = 0.203$ inside $37\%$ of the volume needs a peak of
$1.58\times10^{13}$\,G, which is $0.36\,B_c$ -- \emph{below} the
$0.97\,B_c$ the space-filling $K\rho\varpi$ field reaches. Confining the
field does not price the mass out; it buys headroom, because the confined
$\zeta = 1$ profile is flatter across its region than $K\rho\varpi$ is
across the star, and peak field at fixed energy is set by how peaked the
profile is, not only by how much volume it has.

The ratio is a free parameter over a useful range. At the poloidal
amplitude the companion paper adopts, the confined configuration has
$B_{\mathrm{t}}/B_{\mathrm{p}} = 9.7$ in amplitude and
$E_{\mathrm{t}}/E_{\mathrm{p}} = 884$ in energy; raising $k_0$ by a decade
brings the amplitudes to parity, and both peaks are then still only
$0.36$ and $0.37$ of $B_c$. A configuration with comparable poloidal and
toroidal fields at $2\,M_\odot$ is not excluded by the Landau limit.

What runs out first is the poloidal side, and this is where the
competition actually lives. At $k_0 = 3\times10^{-12}$ the poloidal peak
reaches $1.11\,B_c$ and the equation of state no longer describes the
star it is being used on. The usable window is an exterior dipole from
$\sim6\times10^{10}$ to $\sim10^{12}$\,G. We had expected the competition
to run through the closed-line volume instead -- more escaping flux, a
larger $u_s$, a smaller closed region -- and it does not: that volume
holds at $37.2\%$ across the whole sweep, because a constant $M'(u) = k_0$
makes the Grad--Shafranov source linear in $k_0$, so $u$ and $u_s$ scale
together and the geometry is invariant.

\subsection{What this does not show}

It is a scoping calculation, on a fixed density, and it establishes only
that the energy budget and the field strengths are mutually compatible.
It does not exhibit an equilibrium. The argument that the mass survives
rests on the virial relation containing $E_{\mathrm{mag}}$ as a whole
rather than its distribution, so that a redistribution of the same
magnetic energy leaves the global balance unchanged -- which is sound to
first order and is exactly what the next section has to verify rather
than assume.

%-------------------------------------------------------------------

\section{The non-barotropic construction}
\label{sec:nonbarotropic}

\subsection{Equilibrium without a barotrope}

Dropping $\rho = \rho(P)$ removes the requirement that
$\rho^{-1}\mathbf{f}_L$ be a gradient, but not the requirement that the
Lorentz force be parallel to $\nabla u$, which follows from
Eq.~(\ref{eq:decomposition}) alone. We may therefore write
\begin{equation}
\mathbf{f}_L = g\,\nabla u ,
\qquad
g = -\frac{\Delta^\ast u + \beta\beta'}{4\pi\varpi^2} ,
\label{eq:fl}
\end{equation}
where $g$ is now a function of position and not of $u$ alone. The Euler
equation becomes
\begin{equation}
\nabla P + \rho\nabla\Phi = g\,\nabla u .
\label{eq:euler}
\end{equation}
Taking the component along a surface of constant $\Phi$ gives
\begin{equation}
\mathrm{d}P = g\,\mathrm{d}u
\qquad\text{along equipotentials,}
\label{eq:dPdu}
\end{equation}
which determines $P$ everywhere once it is fixed on one curve crossing
every equipotential; we integrate outward along equipotentials from the
polar axis, where we impose
$P(\text{axis}) = \alpha\,P_{\mathrm{ff}}(\text{axis})$ with
$P_{\mathrm{ff}}$ the pressure of the field-free star at the same central
density. The dimensionless anchor $\alpha$ is a free parameter of the
construction and, as Sect.~\ref{sec:results-rot} shows, it is the knob
that controls the composition the solution demands. The radial component
of Eq.~(\ref{eq:euler}) then returns the density,
\begin{equation}
\rho = \frac{g\,\partial_r u - \partial_r P}{\partial_r \Phi} ,
\label{eq:rhoinv}
\end{equation}
and $\Phi$ is recovered from $\rho$ by solving Poisson's equation. The
three steps are iterated to convergence.

The logic is worth stating plainly, because it is what the rest of the
paper tests. The poloidal source is now free: we may choose
$\Delta^\ast u$ as we like, and the pressure adjusts. Nothing is
determined by an equation of state, because no equation of state has been
imposed. What the solution owes us instead is a demonstration that the
$(P,\rho)$ pair it produces corresponds to matter that could exist.

\subsection{The composition closure, and what it can and cannot absorb}
\label{sec:closure}

We close the system with a zero-temperature degenerate electron gas of
variable composition. Then $P = A f(x)$ with $A$ independent of the mean
molecular weight per electron, while $\rho = B(\mu_e)\,x^3$ with
$B \propto \mu_e$. The relativity parameter $x$ is therefore fixed by
$P$ alone, and
\begin{equation}
\mu_e \;\propto\; \frac{\rho}{x^3(P)} ,
\label{eq:mue}
\end{equation}
so that $\mu_e$ is not an input but a diagnostic: it is read off
pointwise as the mismatch between the density and the pressure the
solution carries. Less pressure supporting a given density means a larger
$\mu_e$, and more pressure means a smaller one.

This is the precise sense in which non-barotropy is composition
stratification here. Fixing $\mu_e$ at a single value everywhere would
restore $P = P(\rho)$ and return us to Sect.~\ref{sec:ceiling}. The
freedom that lets the construction exceed the barotropic ceiling is
exactly the freedom in $\mu_e$, and the question the construction must
answer is whether the $\mu_e$ field it asks for is one a white dwarf can
supply.

The two directions are not symmetric, and this asymmetry turns out to
decide the result.
\begin{itemize}
\item $\mu_e > 2$ is ordinary. Neutron-rich species supply it:
      $^{56}$Fe gives $2.154$, $^{22}$Ne gives $2.2$. A cold interior can
      be built with them.
\item $\mu_e < 2$ has no cold realisation at all. Every $\alpha$-chain
      nucleus --- $^{4}$He, $^{12}$C, $^{16}$O, $^{20}$Ne, $^{24}$Mg ---
      has $A/Z = 2$ exactly, and going below requires hydrogen, which is
      not available in the interior of an ultramassive white dwarf.
\end{itemize}
A solution demanding $\mu_e < 2$ is therefore not asking for a
composition. By Eq.~(\ref{eq:mue}) it is asking for more pressure than
cold degenerate matter provides at that density, and the only physical
source of that excess is finite temperature. We return to this in
Sect.~\ref{sec:thermal}, where it becomes the decisive constraint.

We accept $\mu_e \in [2.0, 2.15]$ as the strict window and report
$[2.0, 2.2]$ separately, since $^{22}$Ne is a defensible contaminant of a
C/O interior.

\subsection{The toroidal component}
\label{sec:toroidal}

Equation~(\ref{eq:beta}) and the confinement to $u > u_s$ survive the
loss of barotropy, as Sect.~\ref{sec:forced} established, so we keep the
form of Eq.~(\ref{eq:betaform}) and normalise $\beta_0$ to a target
$E_{\mathrm{tor}}/|W|$. Two details of the implementation are not
cosmetic, and we record them because each one produced a converged but
spurious answer before it was fixed.

First, $u_s$ must be recomputed against the current star at every outer
iteration, not held at the field-free background. Freezing it left the
toroidal field confined to the region belonging to a $1.74\,M_\odot$
configuration while the iteration converged on a $1.21\,M_\odot$ one, so
part of the field lay outside the matter with nothing to confine it. The
resulting virial error was $0.17$ and was stable from $32$ to $256$ outer
iterations: converged, and converged on something that was not an
equilibrium.

Second, $\zeta = 2$ rather than $\zeta = 1$. At $\zeta = 1$ the
derivative $\beta'$ jumps across the separatrix, which is a current
sheet. Localising the residual of Eq.~(\ref{eq:euler}) showed it piled up
in the equatorial plane at $0.4$--$0.6\,R_{\mathrm{eq}}$ --- on the
separatrix --- and carried essentially the whole virial error; raising
$\zeta$ to $2$, so that $\beta$ and $\beta'$ both vanish there, cut the
error at $E_{\mathrm{tor}}/|W| = 0.10$ from $0.029$ to $0.010$. This is
the same choice made by \citet{ciolfi2013} and \citet{lander2009}, and
for the same reason.

\subsection{Differential rotation}
\label{sec:rotation}

We add rotation on the $j$-constant law,
\begin{equation}
\Omega(\varpi) = \Omega_c\,\frac{A^2}{A^2 + \varpi^2} ,
\qquad
\Psi(\varpi) = -\frac{\Omega_c^2 A^2}{2}\,
               \frac{\varpi^2}{A^2 + \varpi^2} ,
\label{eq:rot}
\end{equation}
with $A$ quoted as a fraction of the background equatorial radius;
$A \to \infty$ is uniform rotation. Because $\Omega$ depends on $\varpi$
alone the centrifugal force is conservative, and the whole construction
goes through with $\Phi$ replaced by $\Phi_{\mathrm{eff}} = \Phi + \Psi$
in Eqs.~(\ref{eq:dPdu}) and (\ref{eq:rhoinv}). Gravity itself is still
solved from $\rho$ alone. We normalise $\Omega_c$ to a target
$T/|W|$, renormalising against the current $|W|$ at each iteration so
that the quoted ratio is the one attained at the solution and not at the
starting guess.

%-------------------------------------------------------------------

\section{Numerical method, and what the acceptance tests cost}
\label{sec:method}

The meridional grid is $129 \times 129$ in $(r,\theta)$ with Poisson and
Grad--Shafranov solved to $l_{\max} = 16$, and $32$ outer iterations;
doubling either the grid or the iteration count moves the reported
virial errors by less than one part in $10^{2}$ of their own value.

A configuration is accepted only if it passes, simultaneously:

\begin{description}
\item[Equilibrium.] The virial error
      $\mathrm{VE} = |2T + W + 3\!\int\!P\,\mathrm{d}V +
      E_{\mathrm{mag}}|/|W| < 10^{-3}$. This is not a formality. We
      verified each term separately against its own identity ---
      $\int \mathbf{x}\cdot\rho\nabla\Phi\,\mathrm{d}V = -W$ to
      $1$ part in $10^4$,
      $\int \mathbf{x}\cdot\nabla P\,\mathrm{d}V = -3\!\int\!P\,\mathrm{d}V$
      to $7$ parts in $10^4$, and
      $\int \mathbf{x}\cdot\mathbf{f}_L\,\mathrm{d}V = E_{\mathrm{mag}}$
      to $1.3$ parts in $10^3$ --- and confirmed that VE equals the
      integrated radial force imbalance. A configuration with
      $\mathrm{VE} = 0.26$ has, in the worst $5\%$ of its volume, a
      pointwise force residual of $0.35$ of the local gravity, against
      $1.7\times10^{-4}$ for an accepted one.
\item[Composition.] Mass-weighted percentiles $\mu_e^{(1)}$ and
      $\mu_e^{(99)}$ inside the window, and at most $1\%$ of the mass
      outside it. The gate must be pointwise. A shell-averaged version
      reported $\mu_e \in [2.048, 2.630]$ for a configuration whose
      pointwise maximum was $65.2$, and reported $[2.002, 2.122]$ ---
      inside the strict window --- for one whose pointwise maximum was
      $2.627$. Mass weighting is what handles the stellar surface, where
      the inversion of Eq.~(\ref{eq:mue}) is ill-conditioned as
      $P \to 0$: that shell holds almost no mass and cannot move a
      mass-weighted percentile, so no \textit{ad hoc} radial cut is
      needed.
\item[Stratification.] The composition profile must be measurable above
      the solver's own noise, and Ledoux-stable. We calibrate the noise
      by running the construction at $k_0 = 0$, where the answer must be
      homogeneous; the residual spread is $0.0203$, and we require a
      profile at least $3\times$ that before computing a Ledoux
      criterion at all. Profiles below that threshold are not declared
      stable, they are declared unmeasured.
\item[Landau.] $\max|\mathbf{B}| < B_c = 4.414\times10^{13}$\,G, applied
      to the total field.
\item[Core support.] The mass-weighted radial Lorentz force inside
      $0.2\,R$, as a fraction of local gravity, must be positive.
\item[Rotation.] No mass shedding, and the attained $T/|W|$ within $5\%$
      of its target. The second test matters: the iteration runs away at
      large $T/|W|$ with concentrated shear, and without a guard it
      returned $M = 2\times10^{13}\,M_\odot$ with $\mathrm{VE} = 1705$ as
      though it were a measurement. Of the $780$ rotating models
      attempted, $334$ were rejected this way.
\end{description}

%-------------------------------------------------------------------

\section{Results without rotation}
\label{sec:results-mixed}

We scanned $128$ non-rotating models over
$\rho_c \in \{10^9, 3\times10^9\}$\,g\,cm$^{-3}$,
$k_0 \in [1.0,2.5]\times10^{-12}$,
$\alpha \in [0.70, 1.00]$ and
$E_{\mathrm{tor}}/|W| \in \{0, 0.005, 0.01, 0.02\}$.

\emph{Not one is admissible.} The individual pass rates are
$104/128$ on equilibrium, $48/128$ on the Landau limit,
$128/128$ on core support, and $2/128$ on composition; the intersection
is empty.

Three results are worth separating out.

\paragraph{The toroidal field does not add mass here.} Along any family
at fixed $(k_0,\alpha,\rho_c)$ the mass falls monotonically as the
toroidal component is turned up: $1.407$, $1.386$, $1.366$ and
$1.279\,M_\odot$ at $E_{\mathrm{tor}}/|W| = 0$, $0.005$, $0.01$ and
$0.02$. What rises with it is $\mu_e$, and the rise is localised: the
maximum of $\mu_e$ migrates from $0.98\,R$ --- a surface artefact of the
ill-conditioned inversion --- to $0.49\,R$ in the equatorial plane, which
is where $|B_\varphi|$ peaks, at $0.48\,R$. The toroidal field takes over
part of the support locally, and Eq.~(\ref{eq:mue}) records that as a
local demand for heavier composition.

\paragraph{Equilibrium itself limits the toroidal fraction.} The virial
error crosses $10^{-3}$ between $E_{\mathrm{tor}}/|W| = 0.01$ and
$0.02$. In consequence $B_\varphi/|\mathbf{B}_{\mathrm{p}}|$ never
exceeds about $0.6$ in any model that is in equilibrium. The
toroidal-dominated interior is not reached.

\paragraph{The stratification is Ledoux-unstable.} Of the $126$ models
whose composition profile rises above the noise floor, \emph{every one}
is Ledoux-unstable, $43$ of them maximally so. The two models that pass
the Ledoux test do so only because their stratification falls below the
threshold and is therefore unmeasured, and both of those fail
composition. The two models with admissible composition are both purely
poloidal, at $1.262$ and $1.336\,M_\odot$, and both are maximally
Ledoux-unstable.

This last point is the mechanism, and it is not a numerical accident.
The extra support that non-barotropy buys is bought by a composition
gradient, and the gradient it needs has the heavy material on the
outside.

%-------------------------------------------------------------------

\section{Results with differential rotation}
\label{sec:results-rot}

We then scanned $780$ models adding
$T/|W| \in \{0, 0.02, 0.04, 0.06, 0.08\}$ and
$A/R_{\mathrm{eq}} \in \{0.5, 1.0, 2.0\}$, and extending the pressure
anchor down to $\alpha = 0.40$; $334$ ran away and were rejected, leaving
$446$.

\paragraph{Mass stops being the difficulty.} Rotation delivers it. The
heaviest model in genuine equilibrium reaches $2.851\,M_\odot$ at
$T/|W| = 0.080$, with $\mathrm{VE} = 2.4\times10^{-4}$, $B < 0.56\,B_c$,
no mass shedding, and its core held. This is a different situation from
the barotropic toroidal branch, which reached $2\,M_\odot$ and then
collapsed dynamically within a few dynamical times.

\paragraph{Mass and composition trade against each other, monotonically.}
Table~\ref{tab:trade} gives the median $\mu_e^{(1)}$ against mass over
the models that are in equilibrium, below $B_c$, on their rotation
target and not shedding. There is no elbow and no island: every increment
of mass is paid for in $\mu_e$, and the payment runs the wrong way ---
towards $\mu_e < 2$, the direction Sect.~\ref{sec:closure} showed has no
cold realisation.

\begin{table}
\caption{Mass against the composition it demands, over the $213$
         rotating models in equilibrium, below $B_c$, on their rotation
         target and not shedding. $\mu_e^{(1)}$ is the mass-weighted
         first percentile; the last column is the pressure excess over
         cold $\mu_e = 2$ matter that the median value implies.}
\label{tab:trade}
\centering
\begin{tabular}{lccc}
\hline\hline
$M$ ($M_\odot$) & $n$ & median $\mu_e^{(1)}$ & implied $\Delta P/P$ \\
\hline
$< 1.3$      & 40 & 2.045 & --- \\
$1.3$--$1.5$ & 48 & 1.912 & $6\%$ \\
$1.5$--$1.8$ & 61 & 1.780 & $16\%$ \\
$1.8$--$2.2$ & 38 & 1.549 & $42\%$ \\
$2.2$--$3.0$ & 26 & 1.363 & $62\%$ \\
\hline
\end{tabular}
\end{table}

\paragraph{Rotation does not lift the toroidal ceiling.} At
$T/|W| = 0.04$ the virial error runs $1.2\times10^{-4}$, $3.2\times
10^{-3}$, $1.7\times10^{-2}$ and $0.26$ at
$E_{\mathrm{tor}}/|W| = 0.005$, $0.02$, $0.05$ and $0.10$. The gate is
crossed in the same place as without rotation; the ceiling moves from
about $0.01$ to about $0.015$, which is not a qualitative change. The
$2.10\,M_\odot$ model that appears at $E_{\mathrm{tor}}/|W| = 0.10$ has
$\mathrm{VE} = 0.26$ and is not an equilibrium.

\paragraph{One model is admissible, and it is unremarkable.} Nothing
passes every test with the strict window. Widening it to
$\mu_e \le 2.2$, which $^{22}$Ne justifies, admits exactly one model out
of $780$:
\begin{quote}
$M = 1.278\,M_\odot$, $\rho_c = 3\times10^9$\,g\,cm$^{-3}$,
$\alpha = 0.55$, $T/|W| = 0.020$, $A = R_{\mathrm{eq}}$,
$B_\varphi/|\mathbf{B}_{\mathrm{p}}| = 0$,
$\mu_e \in [2.020, 2.196]$ with $0.5\%$ of the mass outside,
Ledoux $0.259$, $E_{\mathrm{mag}}/|W| = 0.011$, $B = 0.81\,B_c$,
$\mathrm{VE} = 5.3\times10^{-4}$.
\end{quote}
It is the first fully certified non-barotropic configuration we have
obtained, and it carries no toroidal field, sits below the Chandrasekhar
mass, and --- the point that should not be buried --- is \emph{lighter}
than its own field-free background, $1.278$ against $1.387\,M_\odot$.
The construction that was meant to add support subtracted $8\%$ of the
mass. Since $\alpha < 1$ lowers the central pressure, this is
self-consistent; but it means that not one admissible model in this
paper is held up by its field.

We note that the $1\%$ mass-fraction test was evaluated against the
strict window even when the percentile window was widened, so it is
stricter than the widened window and the count of one is a lower bound.

%-------------------------------------------------------------------

\section{What the composition demand means physically}
\label{sec:thermal}

Every high-mass model asks for $\mu_e$ well below $2$, and
Sect.~\ref{sec:closure} showed that this is a request for pressure, not
for composition. In the ultrarelativistic limit $P \propto
\mu_e^{-4/3}$, so the excess required is
$(2/\mu_e)^{4/3} - 1$. Supplying it as ideal ion thermal pressure of
$^{12}$C, $P_{\mathrm{ion}} = \rho k T / (A m_u)$, fixes the temperature
in Table~\ref{tab:temp}.

\begin{table}
\caption{Temperature that would have to supply the missing pressure, at
         $\rho_c = 3\times10^9$\,g\,cm$^{-3}$ where
         $P_{\mathrm{deg}} = 2.12\times10^{27}$\,erg\,cm$^{-3}$.}
\label{tab:temp}
\centering
\begin{tabular}{lcc}
\hline\hline
$\mu_e$ demanded & $\Delta P / P$ & $T$ (K) \\
\hline
1.945 & $3.8\%$  & $3.9\times10^{9}$ \\
1.90  & $7.1\%$  & $7.2\times10^{9}$ \\
1.70  & $24\%$   & $2.5\times10^{10}$ \\
1.50  & $47\%$   & $4.8\times10^{10}$ \\
\hline
\end{tabular}
\end{table}

At $\rho_c = 3\times10^9$\,g\,cm$^{-3}$ carbon is already in the
pycnonuclear regime, where C+C proceeds without thermal assistance at
all. Requiring in addition a temperature of several times $10^9$\,K puts
the configuration an order of magnitude above any carbon ignition
threshold, and the $1.8$--$2.2\,M_\odot$ band of Table~\ref{tab:trade}
requires $2$--$5\times10^{10}$\,K, which is photodisintegration and NSE
territory rather than white dwarf territory.

The conclusion is not that the thermal route is unexplored but that it is
self-defeating: a configuration in this mass range, supported this way,
would have ignited before it could be assembled. We emphasise that this
is an order-of-magnitude argument with an ideal ion gas and not a solved
finite-temperature equation of state; radiation pressure is subdominant
at these densities. The margin, however, is a factor of $5$ to $100$
above ignition, so the conclusion does not rest on the approximation.

It is worth adding that every realistic refinement of the equation of
state moves in the same, unhelpful direction. Coulomb and lattice
corrections lower the pressure by one to two percent here; general
relativity lowers the maximum mass; Landau quantisation softens the
equation of state. None of them supplies the $24\%$ that
$1.8\,M_\odot$ needs. The Landau limit, incidentally, is not what
closes the route: the certified model sits at $0.81\,B_c$, comfortably
below.

%-------------------------------------------------------------------

\section{Conclusions}
\label{sec:conclusions}

We set out to determine whether abandoning barotropy allows a white dwarf
to carry a strong mixed poloidal--toroidal field, and, with differential
rotation added, to reach the $1.8$--$2\,M_\odot$ range of interest for
ultramassive white dwarfs. It does not, and the obstruction is specific
enough to be useful.

\begin{enumerate}
\item Measured in this code, the barotropic twisted torus saturates at
      $E_{\mathrm{tor}}/E_{\mathrm{mag}} = 0.04$ and
      $E_{\mathrm{pol}}/|W| = 0.018$ (Sect.~\ref{sec:ceiling}).
\item The non-barotropic construction can exceed that, but the
      composition stratification it requires to do so is Ledoux-unstable
      in every model where the stratification is large enough to be
      measured at all --- $126$ of $126$ (Sect.~\ref{sec:results-mixed}).
\item Equilibrium, not composition, limits the toroidal fraction:
      the virial error crosses $10^{-3}$ near
      $E_{\mathrm{tor}}/|W| \simeq 0.01$, and
      $B_\varphi/|\mathbf{B}_{\mathrm{p}}| \lesssim 0.6$ in every model
      that is in equilibrium, with or without rotation. A
      toroidal-dominated interior was not obtained.
\item With differential rotation, mass ceases to be the difficulty ---
      $2.85\,M_\odot$ in genuine equilibrium --- but mass and composition
      trade monotonically, and above $1.5\,M_\odot$ the required $\mu_e$
      falls below $2$, which no cold composition supplies
      (Table~\ref{tab:trade}).
\item Supplying that deficit thermally requires $4\times10^9$\,K at
      $1.4\,M_\odot$ and $2$--$5\times10^{10}$\,K at $2\,M_\odot$,
      against a pycnonuclear carbon ignition that is already active at
      these densities (Sect.~\ref{sec:thermal}).
\item Of $908$ models across both scans, exactly one satisfies every
      test, at $1.278\,M_\odot$, with no toroidal field, and lighter than
      its own field-free background.
\end{enumerate}

The astrophysical reading is that an ultramassive white dwarf held up in
this way does not wait. The support it would need corresponds to
conditions under which carbon has already ignited, which bears directly
on the question of whether such objects end as thermonuclear supernovae
or collapse, and answers it from the side of the equilibrium structure
rather than from the evolution.

What would change the answer is a genuinely self-consistent mixed
Grad--Shafranov solve, in which $u$ responds to $\beta$ instead of being
taken from a purely poloidal solution. We expect that to raise the
toroidal ceiling of point~3, since that ceiling is a failure of force
balance rather than of composition. We do not expect it to help with
points~4 and~5, and it will probably make them worse, since the
composition demand grows with field strength.

\begin{acknowledgements}
TBD.
\end{acknowledgements}

\bibliographystyle{aa}
\bibliography{paper}

\end{document}
